% SGA 3, Expose VI_A, Loop 1 English reconstruction.
% Scope: paragraphs 4.5 and 4.6 in full, through Proposition 4.6.1;
% section 5 is excluded.
% Authority: Polo--Gille Exp6A-13oct24.pdf, component-local page 29
% (combined-reader page 332; printed page 321).
% OCR and the Reinhold English Markdown are comparison witnesses only.

\medskip
\phantomsection\label{par:sga3-VIA-4.5}%
\noindent\textbf{4.5.---}
至此已构造出 \(F\backslash G\), 并证明它是 \(A\) 上有限型概形的直和.
\hyperref[thm:sga3-VIA-3.2]{定理~3.2} 的其余断言直接化归为关于群胚
\(\mathscr S(G^\alpha_*)\) 的断言.
与第~V~篇 \S~6 一样, \SGARefLink{thm:sga3-VIA-3.2:item:i}{(i)} 中的第二个断言
由第一个断言结合 \SGARefLink{thm:sga3-VIA-3.2:item:ii}{(ii)} 与
\SGARefLink{thm:sga3-VIA-3.2:item:iii}{(iii)} 得到;
因此只需证明 \SGARefLink{thm:sga3-VIA-3.2:item:ii}{(ii)}, (iii) 和 (iv)(a).

由于 \(A'\) 是在 \(A\) 上有限且自由的局部 \(A\)-代数,
态射 \(A\to A'\) 忠实平坦且有限呈示.
由 SGA~1, VIII (3.1, \SGARefLink{par:sga3-VIA-4.6}{4.6}, 5.4),
只需对群胚
\[
  \mathscr S(G^\alpha_*)\otimes_AA'.
\]
验证相应断言. 这个群胚同构于有限个 \(\mathscr S(G^0_*)\otimes_AA'\) 的副本的直和
(参见 \hyperref[par:sga3-VIA-4.4]{4.4}), 因而化归到群胚 \(\mathscr S(G^0_*)\).
对于这个群胚, 我们继续按第~V~篇 \S~6 给出的证明进行,
这也是我们在 \SGARefLink{par:sga3-VIA-4.3}{4.3} 中已经开始采用的证明.

\medskip
\phantomsection\label{par:sga3-VIA-4.6}%
\noindent\textbf{4.6.---}
本节最后对第~V~篇的\SGARefLink{lem:sga3-VIA-6.1}{引理~6.1}
和 \S~9(a) 作一些说明.
在第~V~篇 \S~9(a) 的假设与记号下,
我们寻找一个条件, 使得 \((\mathrm{Sch}/S)\)-群胚 \(X_*\) 的余核的形成
与基扩张
\[
  \pi:S'\longrightarrow S.
\]
可交换. \(X_*\) 和 \(X'_*\) 的余核分别等同于
\(X_*\) 和 \(X'_*\) 在拟截面 \(U\) 和 \(U'\) 上诱导的群胚
\(U_*\) 和 \(U'_*\) 的余核.
因此, 化归到满足\SGARefLink{thm:sga3-V-4.1}{第~V~篇定理~4.1} 的假设的
\((\mathrm{Sch}/S)\)-群胚的情形.

\footnote{\label{note:sga3-VIA-45}编者注: 以下这句话为增补内容.}%
设 \(Y\) 为 \(U_*\) 的余核, 置
\[
  Y'=Y\times_SS',
\]
并设 \(Y_1\) 为 \(U'_*\) 的余核;
我们在第~V~篇 \S~9(a) 中已看到, 典范态射
\[
  Y_1\longrightarrow Y'
\]
是同胚 (事实上是普遍同胚); 因此可将 \(Y_1\) 和 \(Y'\) 作为拓扑空间等同.
若 \(p:U\to Y\) 为典范态射,
\(p':U'\to Y'\) 由它经基变换得到,
我们希望以下 \(\mathcal O_{Y'}\)-模序列正合:
\[
  \tag{*}\label{eq:sga3-VIA-4.6-star}
  \mathcal O_{Y'}
    \longrightarrow p'_*(\mathcal O_{U'})
    \rightrightarrows
  p'_*u'_{1*}(\mathcal O_{U'_1})
    =
  p'_*u'_{0*}(\mathcal O_{U'_1}).
\]
\footnote{\label{note:sga3-VIA-46}编者注: 下文对原文作了修改;
原文此处对 \(X_*\) 施加的附加假设并无必要.}%
由于我们处在\SGARefLink{thm:sga3-V-4.1}{第~V~篇定理~4.1} 的假设下,
\(u_0\) 和 \(u_1\) 有限且局部自由; 此外, 由
\SGARefLink{thm:sga3-V-4.1:item:ii}{第~V~篇 4.1(ii)}, \(p\) 为整态射.
所以 \(p\) 和 \(p\circ u_i\) 仿射, 因而分离且拟紧.

因此, 若 \(S'\) 在 \(S\) 上平坦,
EGA~III\textsubscript{1}, 1.4.15 (并计入 EGA~III\textsubscript{2}
中的勘误 Err\textsubscript{III}~25) 将序列
\eqref{eq:sga3-VIA-4.6-star} 等同于下述序列的逆像:
\[
  \tag{**}\label{eq:sga3-VIA-4.6-double-star}
  \mathcal O_Y
    \longrightarrow p_*(\mathcal O_U)
    \rightrightarrows
  p_*u_{1*}(\mathcal O_{U_1})
    =
  p_*u_{0*}(\mathcal O_{U_1}),
\]
而后者是正合的.%
\footnote{\label{note:sga3-VIA-47}编者注: 以下这句话及下面的编号
\SGARefLink{prop:sga3-VIA-4.6.1}{4.6.1} 为增补内容, 以突出所述结果.}
当群胚 \(X_*\) “局部地”具有拟截面时,
也可采用类似论证 (参见\SGARefLink{thm:sga3-V-7.1}{第~V~篇定理~7.1} 的证明).
因此得到:

\medskip
\phantomsection\label{prop:sga3-VIA-4.6.1}%
\noindent\textbf{命题 4.6.1.}
\textit{当 \(X_*\) 局部地具有拟截面时,
\(X_*\) 的余核的形成与平坦基扩张可交换.}
